% Policy Memorandum and Concept Paper. Built by the future pass from the
% instructions below. Do not process now. This is the policy-development-before-
% drafting section: concept paper, legislative intent, and the public-document
% acceleration argument. Keep file names in the table, not the prose.

[Open the section as a formal policy memorandum and concept paper, written as a
top legislative author would: a short memo header (to the relevant committees of
jurisdiction, from the sponsoring office, re the VVUQ Physical AI Oncology Trial
Bill), then connected prose. Do not narrate the source files; synthesise them.
All source files for this section are listed in the table below; abbreviate leaf
names in the prose and cite the works by their bibliography keys.]

[Subsection 1, Concept and legislative intent. State the concept in the prompt
thesis terms: a United States bill that elevates the VVUQ code-verification
process above the generated code itself and requires verification ahead of
generation for any robot-patient interaction in a Physical AI oncology trial.
Ground the framing in the converging fronts argument from the VVUQ-01
introduction (in silico simulation, autonomous code generation, robotic
procedures) and the asymmetry of cost between a pre-delivery error caught at a
gate and a post-delivery failure in a patient. Use the Prefatory Note of the
template (\texttt{template-paper/chunk\_01\_front\_matter.md}) as the model for
tone and for the USL, PSL, and national MCP-PAI vocabulary introduced here and
carried throughout. Cite \texttt{kawchak2026vvuq01}, \texttt{kawchak2026national},
and \texttt{kawchak2026bills}.]

[Subsection 2, Why a law and why now. Argue, from the national-platform executive
summary, that existing frameworks can be adapted rather than replaced, that
Physical AI complexity now exceeds what manual review accommodates at scale, and
that the 2026 FDA real-time clinical trials initiative makes faster automated
assurance timely. Make the patient-safety-and-efficacy purpose explicit and keep
mentions of the FDA and other bodies respectful and opportunistic for advancing
the field and keeping the United States first in surgical-humanoid VVUQ. Cite
\texttt{fda2026realtime}, \texttt{fda2023credibility}, and \texttt{iche6r3}.]

[Subsection 3, Repository-based AI as a legislative accelerator. Develop the
second half of the thesis: state-of-the-art repository-based AI models synthesise
VVUQ code generation and execution results into public, reproducible documents
(this very bill is one), shortening the path from technical evidence to statutory
text. Draw the synthesis pattern from the national-platform source-documents
catalog (how adapted standards and open repositories become legislative
blueprints) and from the VVUQ-02 lineage (instruction to codegen to execution to
paper). Cite \texttt{repo-cancer-automated}, \texttt{repo-physical-ai}, and
\texttt{anthropicclaudecode}.]

[Build one source-file table with the format below; give the parent paths once at
the top rows and abbreviate the leaf names in the body. Reference the table from
the prose; do not list these paths inline.]

% Model table (left-aligned, ragged-right; longer parent path at the top):
% \begin{table}[h]
% \centering
% \begin{tabularx}{\textwidth}{L{4.7cm} Y}
% \toprule
% \textbf{Source file (abbreviated leaf)} & \textbf{What to take from it} \\
% \midrule
% \multicolumn{2}{l}{\textit{Parent: papers/VVUQ-03/template-paper/}} \\
% chunk\_01\_front\_matter.md & Prefatory-note tone; USL, PSL, MCP-PAI vocabulary \\
% \multicolumn{2}{l}{\textit{Parent: papers/VVUQ-01/final-paper/sections/}} \\
% introduction.tex & Converging fronts; cost asymmetry; verification-first thesis \\
% \multicolumn{2}{l}{\textit{Parent: physical-ai-oncology-trials/national-platform/new\_paper/final\_paper/sections/}} \\
% executive\_summary.tex & Adapt-not-replace argument; economic and governance framing \\
% source\_documents.tex & How adapted standards and repositories become legislative blueprints \\
% \bottomrule
% \end{tabularx}
% \caption{Sources for the policy memorandum and concept paper.}
% \end{table}
